% ============================================================
%  Tensegrity Tumbleweed Robot — Key Equations (LaTeX)
%  Copy-paste these into your IEEE paper body as needed.
% ============================================================

% --------------------------------------------------
% 1. DYNAMICS EQUATION
% --------------------------------------------------

% Newton's second law for each node
\begin{equation}
\label{eq:dynamics}
\ddot{\mathbf{q}}_i = \frac{\mathbf{F}_i}{m_{\mathrm{node}}}
\end{equation}

% Total force on node i
\begin{equation}
\label{eq:total_force}
\mathbf{F}_i = \mathbf{F}_{\mathrm{cable},i} + \mathbf{F}_{\mathrm{bar},i}
             + \mathbf{F}_{\mathrm{grav},i} + \mathbf{F}_{\mathrm{contact},i}
             + \mathbf{F}_{\mathrm{fric},i} + \mathbf{F}_{\mathrm{wind},i}
             + \mathbf{F}_{\mathrm{damp},i}
\end{equation}

% Gravity on an inclined slope
\begin{equation}
\label{eq:gravity}
\mathbf{g} = \begin{bmatrix}
g\sin\theta \\ 0 \\ -g\cos\theta
\end{bmatrix}, \qquad
\mathbf{F}_{\mathrm{grav},i} = m_{\mathrm{node}}\,\mathbf{g}
\end{equation}

% Bar rigidity (bilateral penalty spring)
\begin{equation}
\label{eq:bar_force}
\mathbf{F}_{\mathrm{bar}} = k_{\mathrm{bar}}
\frac{L - L_{0,\mathrm{bar}}}{L}\,\mathbf{d}
\end{equation}

% Bar damping
\begin{equation}
\label{eq:bar_damp}
\mathbf{F}_{\mathrm{bar,damp}} = c_{\mathrm{bar}}\,
\bigl(\mathbf{v}_{\mathrm{rel}} \cdot \hat{\mathbf{d}}\bigr)\,\hat{\mathbf{d}}
\end{equation}

% Ground contact — normal force (penalty method)
\begin{equation}
\label{eq:contact_normal}
F_{n,i} = \max\!\bigl(0,\; k_{\mathrm{gnd}}\,\delta_i
          - d_{\mathrm{gnd}}\,\dot{z}_i\bigr)
\end{equation}

% Ground contact — Coulomb + viscous friction
\begin{equation}
\label{eq:friction}
\mathbf{F}_{t,i} = -\mu_c\,F_{n,i}\,
\frac{\mathbf{v}_{t,i}}{\|\mathbf{v}_{t,i}\| + \varepsilon}
\;-\; \mu_v\,\mathbf{v}_{t,i}
\end{equation}

% Structural damping
\begin{equation}
\label{eq:struct_damp}
\mathbf{F}_{\mathrm{damp},i} = -c_{\mathrm{damp}}\,\mathbf{v}_i
\end{equation}


% --------------------------------------------------
% 2. WIND FORCE EQUATION
% --------------------------------------------------

% Relative wind velocity
\begin{equation}
\label{eq:v_rel}
\mathbf{v}_{\mathrm{rel},i} = \mathbf{v}_{\mathrm{wind}} - \mathbf{v}_i
\end{equation}

% Aerodynamic drag on each exposed node
\begin{equation}
\label{eq:wind_force}
\mathbf{F}_{\mathrm{drag},i} = \tfrac{1}{2}\,\rho_{\mathrm{air}}\,C_d\,
A_{\mathrm{node}}\,\|\mathbf{v}_{\mathrm{rel},i}\|\,\mathbf{v}_{\mathrm{rel},i}
\end{equation}

% Exposure condition
\begin{equation}
\label{eq:wind_condition}
\mathbf{F}_{\mathrm{drag},i} \neq \mathbf{0} \quad
\text{iff} \quad z_i > z_{\mathrm{COM}} - \epsilon_z
\end{equation}


% --------------------------------------------------
% 3. CABLE FORCE EQUATION
% --------------------------------------------------

% Unit direction vector along cable
\begin{equation}
\label{eq:cable_dir}
\hat{\mathbf{d}} = \frac{\mathbf{q}_{n_2} - \mathbf{q}_{n_1}}
{\|\mathbf{q}_{n_2} - \mathbf{q}_{n_1}\|}
\end{equation}

% Slack variable
\begin{equation}
\label{eq:slack}
s = L - L_0
\end{equation}

% Softplus activation (smooth tension-only approximation)
\begin{equation}
\label{eq:softplus}
\sigma(s) = \tfrac{1}{2}\!\left(s + \sqrt{s^2 + \varepsilon}\right),
\qquad \varepsilon = 10^{-6}
\end{equation}

% Cable tension magnitude
\begin{equation}
\label{eq:cable_mag}
f_{\mathrm{cable}} = \max\!\bigl(0,\;
k_{\mathrm{cable}}\,\sigma(s) + u_{m,i}\bigr)
\end{equation}

% Cable force vector
\begin{equation}
\label{eq:cable_force}
\mathbf{F}_{\mathrm{cable}} = f_{\mathrm{cable}}\,\hat{\mathbf{d}}
\end{equation}

% Saturation limits
\begin{equation}
\label{eq:cable_sat}
f_{\mathrm{cable}} =
\begin{cases}
0       & \text{if } L < L_{\min} \\[4pt]
u_{\max} & \text{if } L > L_{\max}
\end{cases}
\end{equation}


% --------------------------------------------------
% 4. NMPC COST FUNCTION
% --------------------------------------------------

% State tracking error
\begin{equation}
\label{eq:track_err}
\mathbf{e}_k = \mathbf{x}_k - \mathbf{x}_{\mathrm{ref}}
\end{equation}

% Control increment
\begin{equation}
\label{eq:ctrl_inc}
\Delta\mathbf{u}_k = \mathbf{u}_k - \mathbf{u}_{k-1}
\end{equation}

% Terminal weight
\begin{equation}
\label{eq:term_weight}
w_k =
\begin{cases}
1 & k < N_p \\
3 & k = N_p
\end{cases}
\end{equation}

% Full NMPC cost function
\begin{equation}
\label{eq:nmpc_cost}
J = \sum_{k=1}^{N_p} \left[
  w_k\,\mathbf{e}_k^\top Q\,\mathbf{e}_k
  \;+\; \mathbf{u}_k^\top R\,\mathbf{u}_k
  \;+\; \Delta\mathbf{u}_k^\top R_d\,\Delta\mathbf{u}_k
\right]
\end{equation}

% Subject to constraints
\begin{equation}
\label{eq:nmpc_constraints}
\text{s.t.} \quad
\|\mathbf{u}_k\|_\infty \leq u_{\mathrm{lim}}, \qquad
\|\Delta\mathbf{u}_k\|_\infty \leq \Delta u_{\mathrm{lim}}
\end{equation}

% Simplified prediction model used inside NMPC
\begin{align}
\label{eq:nmpc_predict}
\mathbf{p}_{k+1} &= \mathbf{p}_k + \Delta t\,\mathbf{v}_k \\
\dot{v}_{x} &= \frac{k_{\mathrm{eff}}}{m_{\mathrm{tot}}}\,u_x
              - c_d\,v_x \\
\dot{v}_{y} &= \frac{k_{\mathrm{eff}}}{m_{\mathrm{tot}}}\,u_y
              - c_d\,v_y \\
\dot{v}_{z} &= \frac{k_{\mathrm{eff}}}{m_{\mathrm{tot}}}\,u_z
              - g - c_d\,v_z
\end{align}

% Motor dynamics (2nd-order model)
\begin{equation}
\label{eq:motor}
\ddot{u}_{m,i} = \omega_n^2\,(u_{c,i} - u_{m,i})
               - 2\zeta\,\omega_n\,\dot{u}_{m,i}
\end{equation}


% ============================================================
%  TABLE 1: SIMULATION PARAMETERS
% ============================================================

\begin{table}[htbp]
\centering
\caption{Simulation Parameters}
\label{tab:sim_params}
\begin{tabular}{l l l}
\hline
\textbf{Parameter} & \textbf{Symbol} & \textbf{Value} \\
\hline
\multicolumn{3}{l}{\textit{Structure}} \\
Number of nodes            & $N$              & 12 \\
Number of bars             & $N_b$            & 6 \\
Number of cables           & $N_c$            & 30 \\
Node mass                  & $m_{\mathrm{node}}$ & 0.12\,kg \\
Robot radius               & $r_{\mathrm{robot}}$ & $\approx 0.95$\,m \\
\hline
\multicolumn{3}{l}{\textit{Stiffness \& Damping}} \\
Cable stiffness            & $k_{\mathrm{cable}}$ & 2\,000\,N/m \\
Bar stiffness              & $k_{\mathrm{bar}}$   & 20\,000\,N/m \\
Ground spring              & $k_{\mathrm{gnd}}$   & 5\,000\,N/m \\
Ground dashpot             & $d_{\mathrm{gnd}}$   & 150\,N$\cdot$s/m \\
Structural damping         & $c_{\mathrm{damp}}$  & 0.5\,N$\cdot$s/m \\
Bar damping                & $c_{\mathrm{bar}}$   & 5.0\,N$\cdot$s/m \\
\hline
\multicolumn{3}{l}{\textit{Contact \& Friction}} \\
Coulomb friction coeff.    & $\mu_c$          & 0.3 \\
Viscous friction coeff.    & $\mu_v$          & 0.5 \\
Ground slope               & $\theta$         & 0.03\,rad \\
Rod collision stiffness    & $k_{\mathrm{col}}$ & 800\,N/m \\
Rod radius                 & $r_{\mathrm{rod}}$ & 0.03\,m \\
\hline
\multicolumn{3}{l}{\textit{Motor Model}} \\
Natural frequency          & $\omega_n$       & 25\,rad/s \\
Damping ratio              & $\zeta$          & 0.7 \\
Max actuator force         & $u_{\max}$       & 40\,N \\
Min cable length           & $L_{\min}$       & 0.50\,m \\
Max cable length           & $L_{\max}$       & 1.40\,m \\
Motor PD gains             & $K_p,\;K_d$      & 80,\;12 \\
\hline
\multicolumn{3}{l}{\textit{Aerodynamics}} \\
Wind velocity              & $\mathbf{v}_{\mathrm{wind}}$ & $[8.0,\;1.0,\;0]$\,m/s \\
Drag coefficient           & $C_d$            & 1.0 \\
Air density                & $\rho_{\mathrm{air}}$ & 1.225\,kg/m$^3$ \\
Drag area per node         & $A_{\mathrm{node}}$  & 0.07\,m$^2$ \\
\hline
\multicolumn{3}{l}{\textit{NMPC}} \\
Prediction horizon         & $N_p$            & 6 \\
Control dimension          & $n_u$            & 3 \\
Control limit              & $u_{\mathrm{lim}}$  & 0.15\,m \\
Rate limit                 & $\Delta u_{\mathrm{lim}}$ & 0.06\,m \\
Prediction time step       & $\Delta t_p$     & 0.05\,s \\
NMPC interval              & ---              & every 5 steps \\
\hline
\multicolumn{3}{l}{\textit{Integration}} \\
Time step                  & $\Delta t$       & 0.05\,s \\
Simulation duration        & $T_{\mathrm{sim}}$ & 30.0\,s \\
ODE solver                 & ---              & \texttt{ode15s} (stiff) \\
\hline
\end{tabular}
\end{table}


% ============================================================
%  TABLE 2: MATERIAL PROPERTIES
% ============================================================

\begin{table}[htbp]
\centering
\caption{Material Properties}
\label{tab:materials}
\begin{tabular}{l c c}
\hline
\textbf{Property} & \textbf{Cable (Dyneema SK75)} & \textbf{Bar (Al 6061-T6)} \\
\hline
Material type       & UHMWPE fiber            & Aluminium alloy tube \\
Young's modulus $E$ & 100\,GPa                & 69\,GPa \\
Tensile strength    & 3\,400\,MPa (UTS)       & 310\,MPa (UTS) \\
Yield / allowable   & 900\,MPa                & 276\,MPa \\
Density $\rho$      & 970\,kg/m$^3$           & 2\,700\,kg/m$^3$ \\
Diameter            & 2\,mm                   & 10\,mm OD / 8\,mm ID \\
Cross-section $A$   & 3.14\,mm$^2$            & 28.27\,mm$^2$ \\
Prestrain           & 0.5\%                   & --- \\
Prestress           & 500\,MPa                & --- \\
\hline
\multicolumn{3}{l}{\textit{Mass Breakdown}} \\
\hline
Bar mass (each)     & ---                     & $\rho_b A_b L_b$ \\
Cable mass (each)   & $\rho_c A_c L_c$        & --- \\
Joint mass (each)   & \multicolumn{2}{c}{20\,g (3D-printed)} \\
Total struct. mass  & \multicolumn{2}{c}{$\approx 1.44$\,kg (sim: 12 $\times$ 0.12\,kg)} \\
\hline
\end{tabular}
\end{table}


% ============================================================
%  TABLE 3: PERFORMANCE RESULTS
% ============================================================

\begin{table}[htbp]
\centering
\caption{Simulation Performance Results}
\label{tab:performance}
\begin{tabular}{l l}
\hline
\textbf{Metric} & \textbf{Value} \\
\hline
\multicolumn{2}{l}{\textit{Locomotion}} \\
Simulation time          & 30.0\,s (600 steps) \\
Target position $(x,y)$  & $(2.0,\;0.0)$\,m \\
Final tracking error     & reported at runtime \\
Distance travelled       & reported at runtime \\
Mean speed               & reported at runtime \\
\hline
\multicolumn{2}{l}{\textit{Energy}} \\
Total actuator energy    & $E_{\mathrm{act}}$\,(cumulative) \\
Locomotion efficiency    & $\eta = d_{\mathrm{trav}} / E_{\mathrm{act}}$\,[m/J] \\
Mechanical energy drift  & reported at runtime \\
\hline
\multicolumn{2}{l}{\textit{Structural Integrity}} \\
Max cable stress         & $<900$\,MPa (allowable) \\
Max bar stress           & $<276$\,MPa (yield) \\
Cable safety factor      & $\mathrm{SF}_c = \sigma_{\mathrm{allow}} / \sigma_{\max,c}$ \\
Bar safety factor        & $\mathrm{SF}_b = \sigma_{\mathrm{yield}} / \sigma_{\max,b}$ \\
Cable utilisation         & $(\sigma_{\max,c}/\sigma_{\mathrm{allow}}) \times 100\%$ \\
Bar utilisation           & $(\sigma_{\max,b}/\sigma_{\mathrm{yield}}) \times 100\%$ \\
\hline
\multicolumn{2}{l}{\textit{State Estimation (UKF)}} \\
State dimension          & 9 (pos$_3$, vel$_3$, shape$_3$) \\
Measurement dimension    & 6 \\
Covariance convergence   & $\mathrm{tr}(\mathbf{P})$ monitored \\
\hline
\end{tabular}
\end{table}


% ============================================================
%  TABLE 4: SYSTEM COMPARISON (Cable Materials)
% ============================================================


\begin{table}[htbp]
\centering
\caption{Cable Material Comparison for Tensegrity Tumbleweed Robots}
\label{tab:material_comparison}
\begin{tabular}{l c c c c}
\hline
\textbf{Property} & \textbf{Dyneema SK75} & \textbf{Kevlar 49} & \textbf{Steel Wire} & \textbf{Carbon Fiber} \\
\hline
Young's modulus [GPa]         & 100   & 112   & 200   & 230 \\
Tensile strength [MPa]        & 3\,400 & 3\,600 & 1\,800 & 3\,500 \\
Density [kg/m$^3$]            & 970   & 1\,440 & 7\,800 & 1\,750 \\
Specific strength [kN$\cdot$m/kg] & 3.51  & 2.50  & 0.23  & 2.00 \\
UV resistance                 & Good  & Poor  & Excellent & Fair \\
Creep resistance              & Moderate & Good & Excellent & Good \\
Cost                          & Medium & High & Low   & Very High \\
\hline
\multicolumn{5}{l}{\textit{Suitability for Tensegrity Tumbleweed}} \\
\hline
Weight penalty                & Low   & Medium & Very High & Medium \\
Flexibility                   & High  & Medium & Low   & Low \\
Fatigue life                  & Excellent & Good & Fair & Good \\
\textbf{Overall rating}       & \textbf{Best} & Good & Poor & Good \\
\hline
\end{tabular}
\end{table}


% ============================================================
%  FUTURE WORK SECTION
% ============================================================

\section{Future Work}
\label{sec:future_work}

While the proposed tensegrity tumbleweed robot demonstrates promising results
in simulation, several avenues for future research remain:

\begin{enumerate}

\item \textbf{Hardware Prototype and Experimental Validation.}
The current work relies on a numerical simulation with penalty-based contact
and simplified aerodynamic models. Fabricating a physical TT-6 prototype
with Dyneema SK75 cables and aluminium 6061-T6 bars would enable
validation of the dynamic model, friction coefficients, and NMPC
controller performance under real-world conditions.

\item \textbf{Stochastic Wind Modelling.}
The present simulation employs a constant wind field
$\mathbf{v}_{\mathrm{wind}} = [8.0,\;1.0,\;0]$\,m/s. Incorporating
turbulent wind models---such as the Dryden or Von K\'{a}rm\'{a}n gust
spectra---would improve fidelity and test the robustness of the NMPC
controller under realistic atmospheric disturbances. Additionally,
online wind estimation could be integrated into the UKF state vector.

\item \textbf{Terrain Adaptability.}
The ground model assumes a planar slope with uniform friction. Future
work should investigate locomotion over uneven terrain, including rocky
surfaces, sand, and deformable substrates, with obstacle avoidance
integrated into the NMPC cost function.

\item \textbf{Advanced Control Strategies.}
The simplified prediction model used in the NMPC (reduced 3-DOF
dynamics) limits optimality. Full nonlinear MPC using the tensegrity
dynamics, reinforcement learning for gait discovery, and Model
Predictive Path Integral (MPPI) control are promising alternatives for
improved real-time performance.

\item \textbf{Planetary Exploration Deployment.}
Tensegrity tumbleweed robots are strong candidates for planetary
surface exploration. Adapting the simulation to Martian conditions
($g = 3.72$\,m/s$^2$, $\rho_{\mathrm{air}} = 0.020$\,kg/m$^3$)
and addressing low-temperature effects on Dyneema creep properties
would assess feasibility for extraterrestrial missions.

\item \textbf{Structural Topology Optimisation.}
The icosahedral TT-6 geometry is fixed in this study. Future work could
explore topology optimisation of the cable network, variable-stiffness
actuators using shape memory alloys (SMAs), and scaled-up TT-12 or
TT-24 configurations to improve payload capacity and locomotion range.

\item \textbf{Multi-Robot Coordination.}
Deploying a swarm of tensegrity tumbleweeds for distributed sensing and
area coverage is a natural extension. Decentralised control algorithms
and inter-robot communication protocols would enable coordinated
exploration over large areas.

\item \textbf{Enhanced State Estimation.}
The current UKF operates on a 9-dimensional reduced state.
Extending to full node-level estimation (72 states) via distributed
filtering, incorporating vision-based localisation (camera + IMU
fusion), and adding contact-state detection for terrain classification
would improve autonomy and situational awareness.

\end{enumerate}
